% =============================================================================
\section{Conclusion}
\label{sec:conclusion}
% =============================================================================

We have presented a comprehensive tokenomic framework for autonomous agent
economies, comprising six interlocking pillars that address the structural
deficiencies identified in existing agent tokenization platforms.

\subsection{Summary of Results}
\label{subsec:summary}

The six-pillar framework provides the following:

\begin{enumerate}
  \item \textbf{Quadratic bonding curve with creator vesting.} The affine
    pricing function $P(q) = B + Sq$ provides transparent, mathematically
    predictable price discovery with convex cost growth that rewards early
    participants without excluding later ones
    (Proposition~\ref{prop:monotonicity}, Theorem~\ref{thm:curve-comparison}).
    The 180-day creator vesting eliminates immediate-exit moral hazard and
    serves as a quality signal in the separating equilibrium of
    Theorem~\ref{thm:separating}.

  \item \textbf{Dynamic fees.} The volatility-adaptive fee function
    $f(\sigma) = f_{\text{base}} + f_{\text{dyn}} \cdot \min(\sigma/\sigma_{\max}, 1)$
    generates 23--31\% more cumulative revenue than flat fees
    (Proposition~\ref{prop:dynamic-revenue}, Figure~\ref{fig:fee-revenue})
    while reducing adverse selection during high-volatility periods and
    minimizing friction during low-volatility periods.

  \item \textbf{veCORTEX staking.} The vote-escrowed mechanism, available from
    day one, creates 3--5$\times$ lifetime yield advantage for early stakers
    (Proposition~\ref{prop:day-one}, Figure~\ref{fig:staker-yield}) and
    supports a Nash equilibrium in which rational agents stake rather than sell
    (Proposition~\ref{prop:stake-nash}, Theorem~\ref{thm:folk}).

  \item \textbf{Buyback-and-burn.} Multi-source revenue feeds a weekly buyback
    with a dual cap, establishing conditions under which the circulating supply
    becomes net deflationary at approximately month 18--24
    (Theorem~\ref{thm:deflationary}, Figure~\ref{fig:supply-trajectory}).

  \item \textbf{Composability rewards.} The MCP-based composability framework
    with Composability Score and x402 micropayments incentivizes inter-agent
    integration, producing superlinear network value growth consistent with
    Metcalfe's law (Figure~\ref{fig:network-value}) with sybil resistance
    via the $\sqrt{\text{unique\_callers}}$ term
    (Proposition~\ref{prop:sybil}).

  \item \textbf{Sovereign governance.} The constitutional governance framework
    with AI Senator, veCORTEX-weighted Lawmaker voting, and liquid delegation
    achieves a Pareto-improving equilibrium
    (Theorem~\ref{thm:pareto}) with moderate voting power concentration
    (Gini~$= 0.38$, Figure~\ref{fig:governance}).
\end{enumerate}

\subsection{Key Quantitative Results}
\label{subsec:key-results}

Across 1{,}000 Monte Carlo simulation runs per scenario
(Section~\ref{sec:simulations}):

\begin{itemize}
  \item \textbf{Deflationary crossover:} Month 13 (aggressive), month 18
    (moderate), month 24 (conservative).
  \item \textbf{Staker yield:} 15--40\% APR, declining over the 4-year
    emission schedule, with longer locks yielding approximately 2.8$\times$
    the APR of shorter locks.
  \item \textbf{Network value:} Superlinear growth with $R^2 = 0.94$
    Metcalfe's law fit, driven by composability rewards that increase
    inter-agent connection density by approximately 8$\times$ relative to
    unrewarded baselines.
  \item \textbf{Crash recovery:} 43\% faster TVL recovery compared to
    platforms without vesting and dynamic fees, driven by locked supply and
    reduced low-volatility friction.
  \item \textbf{Revenue premium:} 27\% median cumulative revenue advantage
    of dynamic over flat fees, with the premium most pronounced during
    volatile market periods.
\end{itemize}

\subsection{Limitations}
\label{subsec:limitations}

Several limitations warrant acknowledgment:

\begin{enumerate}
  \item \textbf{Model assumptions.} The game-theoretic results rely on
    assumptions of rationality, common knowledge of the game structure, and
    specific distributional forms (log-normal volatility, Poisson agent
    arrivals) that may not hold in practice.

  \item \textbf{Algorand ecosystem size.} As of this writing, the Algorand
    ecosystem is smaller than the dominant smart contract platforms. The
    network effects that drive PURECORTEX's value proposition require a
    critical mass of agents and users; achieving this on a smaller ecosystem
    may be more challenging than on established platforms.

  \item \textbf{Real-world adoption uncertainty.} The simulation results are
    conditional on the parameter assumptions described in
    Appendix~\ref{app:simulation}. Actual ecosystem growth may deviate
    significantly from these assumptions, particularly in the early stages
    when network effects have not yet materialized.

  \item \textbf{Regulatory uncertainty.} The legal and regulatory treatment of
    AI agent tokens is evolving. Changes in regulatory frameworks could
    affect the viability of certain mechanisms (e.g., creator vesting could
    be classified as a securities lock-up with compliance implications).

  \item \textbf{AI Senator limitations.} The welfare function
    (Equation~\eqref{eq:welfare}) is a simplified proxy for ecosystem health.
    The Senator's effectiveness depends on the accuracy of this proxy and the
    quality of the underlying AI model, both of which may degrade in novel
    situations outside the training distribution.
\end{enumerate}

\subsection{Future Work}
\label{subsec:future-work}

Several extensions merit investigation:

\begin{enumerate}
  \item \textbf{Cross-chain bridging.} Extending the composability framework
    to agents deployed on other blockchains via cross-chain message passing
    protocols, enabling a multi-chain agent economy with unified CORTEX
    settlement.

  \item \textbf{Progressive DAO transition.} Formalizing a transition path
    from the initial governance structure (with protocol team guardrails) to a
    fully autonomous DAO, with quantitative milestones for each transition
    stage (e.g., minimum staker count, governance participation rate,
    proposal approval rate).

  \item \textbf{Advanced MCP tooling.} Developing specialized MCP tool types
    (streaming data feeds, persistent state access, multi-agent workflow
    orchestration) with corresponding fee and reward structures.

  \item \textbf{Formal verification.} Applying formal verification techniques
    (model checking, theorem proving) to the smart contract implementations to
    provide mathematical guarantees beyond the analytical proofs presented
    herein.

  \item \textbf{Empirical validation.} Deploying the framework on Algorand
    TestNet with synthetic agent populations to validate the simulation
    predictions against real blockchain execution dynamics.
\end{enumerate}
